../cictikz/src/cictikz/data/tex/cictikz_preamble.tex

% The two-bit R-2R ladder in full: every leg carries a binary weighted
% current out of -V_REF, and each bit steers its leg either into the
% amplifier's summing node or into ground. The leg current is the same
% either way, which is the whole point of steering rather than switching the
% leg off: the ladder always sees the same resistance.
%
% The rightmost 2R is the termination dac_2r_0 starts from and has no switch;
% it is what makes the ladder repeat.

\begin{circuitikz}[american, thick, transform shape, circuit ee IEC]

  ../cictikz/src/cictikz/data/tex/cictikz_lib.tex
  ../cictikz/src/cictikz/data/tex/rdac_lib.tex

  \def\ybot{0.4}          % the -V_REF rail
  \def\ytop{5.4}          % the summing rail into the amplifier
  \def\xa{0}              % b_1 leg
  \def\xb{3.2}            % b_0 leg
  \def\xc{5.4}            % termination

  % ---- The bottom rail: -V_REF, then R between the two switched legs ----
  \draw (\xa,\ybot) -- ++(0.8,0) \hresistor{$R$} -- ++(0.8,0);
  \draw (\xb,\ybot) -- (\xc,\ybot);
  \draw (\xa,\ybot) -- ++(-1.2,0) coordinate (vrefport);
  \rdacPort{(vrefport)}{north}{$-V_{REF}$}

  % ---- Three 2R legs ----
  \foreach \x in {\xa,\xb,\xc} {
    \draw (\x,\ybot) \vresistor{$2R$} -- ++(0,0.7);
  }

  % ---- The two switched legs steer into the summing rail or into ground ----
  \rdacSpdt{(\xa,{\ybot+2.3})}{\ytop}{$b_1$}{$\overline{b_1}$}
  \rdacSpdt{(\xb,{\ybot+2.3})}{\ytop}{$b_0$}{$\overline{b_0}$}
  \node[anchor=west] at (\xa+1.1,{\ybot+0.8}) {$I_1$};
  \node[anchor=west] at (\xb+1.1,{\ybot+0.8}) {$I_0$};

  % The termination has no switch: it sits on ground for good.
  \draw (\xc,{\ybot+2.3}) -- ++(0,0.65) -- ++(0.9,0) coordinate (tg);
  \rdacVgnd{(tg)}

  \rdacAmp{OA}{(8.2,4.8)}

  % ---- Summing rail into the amplifier ----
  % The rail runs from the b_1 contact across the b_0 one to the - input.
  \draw (\xa+0.55,\ytop) -- (OAinn);
  \fill (\xb+0.55,\ytop) circle (0.07);

  % ---- Feedback resistor over the top ----
  \draw (7.0,\ytop) -- (7.0,7.6) -- ++(0.9,0) \hresistor{$R_F$} -- ++(0.9,0)
        -- (11.6,7.6) -- (11.6,4.8) -- (OAout);
  \fill (7.0,\ytop) circle (0.07);
  \draw (OAout) -- ++(0.9,0);

  % ---- The + input sits on ground ----
  \draw (OAinp) -- ++(-1.2,0) coordinate (og);
  \draw (og) -- ++(0,-0.3) \vground;

\end{circuitikz}
\end{document}
